\documentclass[11pt]{article}

\usepackage[utf8]{inputenc}
\usepackage[T1]{fontenc}
\usepackage{amsmath, amssymb, amsthm}
\usepackage{mathtools}
\usepackage[margin=1in]{geometry}
\usepackage{hyperref}

\newtheorem{theorem}{Theorem}
\newtheorem{lemma}[theorem]{Lemma}
\newtheorem{proposition}[theorem]{Proposition}
\theoremstyle{definition}
\newtheorem{definition}[theorem]{Definition}
\newtheorem{example}[theorem]{Example}
\theoremstyle{remark}
\newtheorem{remark}[theorem]{Remark}

\newcommand{\F}{\mathcal{F}}
\newcommand{\Om}{\Omega}
\newcommand{\R}{\mathbb{R}}
\newcommand{\N}{\mathbb{N}}

\title{Probability Measures}
\author{Math Foundations -- Node 19}
\date{}

\begin{document}
\maketitle

\section{Setting}

Throughout, $\Om$ is a non-empty set (the \emph{sample space}) and
$\F \subseteq 2^{\Om}$ is a $\sigma$-algebra on $\Om$ in the sense of node 18:
$\Om \in \F$, $\F$ is closed under complements, and $\F$ is closed under
countable unions. Elements of $\F$ are called \emph{events}.

\section{Definitions}

\begin{definition}[Probability measure]
A \emph{probability measure} on $(\Om, \F)$ is a set function $P : \F \to [0,1]$
satisfying the three \emph{Kolmogorov axioms}:
\begin{enumerate}
    \item[\textnormal{(K1)}] \emph{Non-negativity:} $P(A) \geq 0$ for every $A \in \F$.
    \item[\textnormal{(K2)}] \emph{Normalisation:} $P(\Om) = 1$.
    \item[\textnormal{(K3)}] \emph{Countable additivity ($\sigma$-additivity):} for every sequence
    $\{A_n\}_{n \geq 1} \subseteq \F$ of pairwise disjoint events,
    \[
        P\!\left( \bigcup_{n=1}^{\infty} A_n \right)
        \;=\; \sum_{n=1}^{\infty} P(A_n).
    \]
\end{enumerate}
\end{definition}

\begin{definition}[Probability space]
A \emph{probability space} is a triple $(\Om, \F, P)$ where $\Om$ is a sample space,
$\F$ is a $\sigma$-algebra on $\Om$, and $P$ is a probability measure on $(\Om, \F)$.
\end{definition}

\begin{remark}
From (K1)--(K3) one immediately obtains $P(\emptyset) = 0$ (take all $A_n = \emptyset$ in (K3);
the sum $\sum P(\emptyset)$ converges only if each term vanishes) and finite additivity
$P(A \sqcup B) = P(A) + P(B)$ by padding with empty sets.
\end{remark}

\section{The inclusion--exclusion identity}

\begin{theorem}[Inclusion--exclusion, two-set form]\label{thm:incl-excl}
Let $(\Om, \F, P)$ be a probability space. For every $A, B \in \F$,
\[
    P(A \cup B) \;=\; P(A) + P(B) - P(A \cap B).
\]
\end{theorem}

\begin{proof}
We decompose $A \cup B$ and $B$ into disjoint pieces and apply finite additivity,
which follows from (K3) as noted above.

\smallskip
\textbf{Step 1: $A \cup B = A \sqcup (B \setminus A)$ disjointly.}
Every $\omega \in A \cup B$ is in $A$ or in $B \setminus A := B \cap A^c$, and these
two sets are disjoint by construction. Both lie in $\F$: $A \in \F$ by hypothesis,
and $B \setminus A = B \cap A^c \in \F$ since $\F$ is closed under complements and
finite intersections. By finite additivity,
\begin{equation}\label{eq:step1}
    P(A \cup B) \;=\; P(A) + P(B \setminus A).
\end{equation}

\smallskip
\textbf{Step 2: $B = (B \cap A) \sqcup (B \setminus A)$ disjointly.}
Every $\omega \in B$ either lies in $A$ (hence in $B \cap A$) or does not (hence in
$B \setminus A$), and these two sets are disjoint. Both lie in $\F$ for the same
reasons as in Step~1. By finite additivity,
\[
    P(B) \;=\; P(B \cap A) + P(B \setminus A),
\]
and rearranging gives
\begin{equation}\label{eq:step2}
    P(B \setminus A) \;=\; P(B) - P(A \cap B).
\end{equation}
(Note $P(B \cap A) = P(A \cap B)$ since intersection is commutative.)

\smallskip
\textbf{Step 3: combine.} Substituting~\eqref{eq:step2} into~\eqref{eq:step1},
\[
    P(A \cup B) \;=\; P(A) + \bigl( P(B) - P(A \cap B) \bigr)
                \;=\; P(A) + P(B) - P(A \cap B). \qedhere
\]
\end{proof}

\begin{remark}
The general $n$-set inclusion--exclusion formula
\[
    P\!\left( \bigcup_{i=1}^{n} A_i \right)
    \;=\; \sum_{k=1}^{n} (-1)^{k+1} \!\!\!
            \sum_{1 \leq i_1 < \cdots < i_k \leq n} \!\!\! P(A_{i_1} \cap \cdots \cap A_{i_k})
\]
follows by induction on $n$, with Theorem~\ref{thm:incl-excl} as the base case.
\end{remark}

\section{An immediate corollary}

\begin{proposition}[Boole's inequality]
For every countable collection $\{A_n\}_{n \geq 1} \subseteq \F$,
\[
    P\!\left( \bigcup_{n=1}^{\infty} A_n \right) \;\leq\; \sum_{n=1}^{\infty} P(A_n).
\]
\end{proposition}

\begin{proof}
Set $B_1 = A_1$ and $B_n = A_n \setminus \bigcup_{k < n} A_k$ for $n \geq 2$.
The $B_n$ lie in $\F$, are pairwise disjoint, and satisfy
$\bigcup_n B_n = \bigcup_n A_n$ with $B_n \subseteq A_n$. By (K3) and monotonicity
(itself a one-line consequence of Theorem~\ref{thm:incl-excl} applied with $B \supseteq A$),
\[
    P\!\left( \bigcup_n A_n \right) = \sum_n P(B_n) \leq \sum_n P(A_n). \qedhere
\]
\end{proof}

\end{document}
